\documentclass{article}
\title{A tiny paper for the absorb workflow}
\begin{document}
\maketitle
\begin{abstract}
We fit a line and prove a lemma about tiny papers.
\end{abstract}
\section{Introduction}\label{sec:intro}
Tiny papers are read in full or not at all.
We claim that a line fits the points in Figure~\ref{fig:plot}.
\section{Method}\label{sec:method}
The model is a line.
\begin{equation}
  y = m x + b
  \label{eq:line}
\end{equation}
\subsection{Fitting}\label{sec:fit}
We minimize the squared error.
\begin{align}
  L &= \sum_i (y_i - m x_i - b)^2 \\
  \hat m &= \arg\min_m L
\end{align}
\begin{algorithm}
\caption{Fit, carefully}\label{alg:fit}
Compute the least-squares solution.
\end{algorithm}
\begin{figure}
\includegraphics{figures/plot.png}
\caption{Points and the fitted line.}\label{fig:plot}
\end{figure}
\section{Results}\label{sec:results}
The fit has an error of 0.1, shown in Table~\ref{tab:numbers}.
\begin{table}
\caption{Errors per run.}\label{tab:numbers}
\begin{tabular}{ll}
run & error \\
first & 0.1
\end{tabular}
\end{table}
\begin{lemma}\label{lem:tinier}
Some tiny papers are tinier.
\end{lemma}
\appendix
\section{Proof of the lemma}\label{app:proof}
By inspection.
\end{document}
